\chapter{The Audio System}
\label{ch:audio-system}

CNA's \cnans{Microsoft::Xna::Framework::Audio} namespace sits on a foundation the project
debated openly before choosing: real XNA runs XACT content through Microsoft's own XACT
engine; FNA reimplements that engine as FAudio. CNA does neither --- the project owner
explicitly rejected an optional-FAudio-backend proposal, and instead built a hand-written
XACT container parser (\texttt{CNA::Internal::Audio::XactParser}) that feeds real,
SDL3\_mixer-based playback. This single decision is the root cause of essentially every
deviation from FNA's own audio behavior catalogued in this chapter.

\section{Why SDL3\_mixer, and what that choice costs}

SDL3\_mixer's own device model hard-floors every audio device to one fixed format, and CNA
leans into that rather than fighting it: its internal mixer always requests 16-bit signed
stereo at 44100~Hz, regardless of what the physical output device's native format actually
is --- a deliberate simplification, since XACT content is itself almost always authored at
44100~Hz anyway. The cost of building on SDL3\_mixer rather than a dedicated 3D-audio engine
shows up consistently: no true per-source 3D audio graph, no auxiliary-send/reverb bus, and a
single shared ``cooked callback'' slot per track used for both filtering and pan-crossfeed at
once, rather than independent processing stages.

\section{SoundEffect and its instances}

\cnaclass{SoundEffect} loads raw PCM/WAV assets and is move-only, specifically so that
disposing the parent cascades correctly to every \cnaclass{SoundEffectInstance} it ever
handed out via \texttt{CreateInstance()}. Its raw-buffer constructors require
\emph{headerless} 16-bit PCM data --- passing an entire WAV file's bytes, RIFF header
included, is flagged directly in the class's own documentation as a common, specific mistake
to avoid. \cnaclass{SoundEffectInstance} exposes volume/pan/pitch/looping/state and a
\texttt{Apply3D(listener, emitter)} method implementing a distance-and-pan
\emph{approximation} rather than full 3D audio --- SDL3\_mixer has no native Doppler
computation, so \texttt{DopplerScale} only ever feeds this approximation, never a true
per-sample Doppler shift. Pitch shifting uses XNA's exact exponential formula,
$2^{\text{pitch}}$, and a real state-variable filter (low-pass, high-pass, or band-pass) is
implemented via an SDL3\_mixer per-track callback matching FAudio's own filter algorithm
precisely. \cnaclass{DynamicSoundEffectInstance} supports both integer and, as a
\texttt{NOXNA} addition with no real-XNA equivalent, floating-point buffer submission; its
\texttt{PendingBufferCount} only decreases once a submitted chunk is genuinely consumed by
playback, matching FNA's own behavior rather than decrementing eagerly at submission time.

\section{The XACT object model: AudioEngine, SoundBank, WaveBank, Cue}

\cnaclass{AudioEngine} parses a compiled \texttt{.xgs} settings file and owns the lifecycle of
every \cnaclass{SoundBank}/\cnaclass{WaveBank}/\cnaclass{Cue} built from it, including
per-category instance-limit enforcement and both global and cue-scoped XACT runtime
variables. \cnaclass{WaveBank} supports both a fully in-memory mode and a real streaming mode
(reading by offset and packet size rather than loading everything up front).
\cnaclass{SoundBank::Dispose()} force-stops every cue it ever produced, matching the real
FACT engine's own \texttt{FACTSoundBank\_Destroy} behavior exactly. \cnaclass{Cue} implements
the full XACT playback state machine (\texttt{Created}/\texttt{Preparing}/\texttt{Prepared}/%
\texttt{Playing}/\texttt{Stopping}/\texttt{Stopped}) with an independent \texttt{paused}
flag layered on top of \texttt{Playing} --- matching FACT's own bitmask-shaped state rather
than a naive single enum --- real fade-in/fade-out ramps, per-tick re-evaluation of
Runtime Parameter Control curves, and all five of XACT's randomized track-variation selection
algorithms (\texttt{Ordered}, \texttt{OrderedFromRandom}, \texttt{Random},
\texttt{RandomNoRepeats}, \texttt{Shuffle}).

\section{A real user-reported bug, and the audit it triggered}

The most recent, and most consequential, audit of this namespace was not routine --- it was
triggered by real user reports of high-pitched, sped-up, distorted, or missing audio in an
actual ported game. The highest-priority finding directly explains the ``missing audio''
class of report: the XNB \cnaclass{SoundEffectReader} only ever accepted a narrow PCM format,
meaning \emph{any} MS-ADPCM-compressed XACT wave bank or XNB sound effect --- a common,
space-saving compression choice for XACT content --- silently failed to decode at all,
because the reader's synthetic WAV wrapper was missing the coefficient table MS-ADPCM
decoding requires. This is now fixed, and the fix was extended to cover 8-bit PCM, IEEE
float, and IMA-ADPCM as well, not just the one format that triggered the investigation.

A second real fix addressed an asymmetry in how \cnaclass{DynamicSoundEffectInstance}
handled integer versus float buffer submission --- a genuine code-level risk, not merely a
missing feature. And because no end-to-end audible-output conformance test existed at all
before this audit, a new deterministic offline rendering harness was built specifically to
measure this class of bug directly: it empirically rules out SDL3\_mixer's own resampler as
the cause of the reported pitch issue, confirming frequency is preserved within 0.1\% across
22050/44100/48000/96000~Hz source material resampled into the fixed 44100~Hz mixer, provided
the source sample rate is declared correctly.

A separate, real stereo-panning bug was found and fixed for \cnaclass{SoundEffectInstance}:
a hard pan used to eliminate the opposite channel entirely rather than crossfeed-blending it
down, corrected for the \texttt{Play}/\texttt{Apply3D}/pan-setter paths --- but the static,
fire-and-forget \texttt{SoundEffect::Play(volume, pitch, pan)} convenience helper still
carries the old, unfixed behavior, tracked separately rather than silently left inconsistent
without a record.

\section{Accepted deviations, stated as such}

Several differences from a hypothetical ``full XACT'' implementation are documented as
deliberate, accepted deviations rather than open bugs. XACT's three instance-limit-behavior
modes (queue, replace-oldest, replace-quietest) all collapse to a single ``evict the oldest
active cue'' behavior --- matching FAudio's own shipped implementation, which makes the same
collapse, not a CNA-specific shortcut. A cue-level instance limit has no category- or
same-cue victim-search filter. Runtime Parameter Control curves targeting a DSP preset are
never evaluated, because no DSP preset system exists at all. \texttt{Apply3D}'s panning is a
single-axis linear approximation, not full multi-speaker 3D diffusion. A parsed per-track
filter can only ever decode as low-pass or high-pass, never band-pass --- itself a real bit-decoding
quirk in FAudio being faithfully replicated, not invented independently. Reverb is a
documented no-op, since SDL3\_mixer provides no auxiliary send/return bus to implement it
against, and there is no HRTF or elevation modeling in the 3D audio approximation at all.

\section{Real memory-safety bugs found and fixed}

Beyond behavioral correctness, the same audit pass found and fixed several genuine
memory-safety issues through a systematic ``what happens if operation X runs concurrently
with operation Y'' review: a mixer generation-counter pair with a real use-after-free,
a heap-buffer-overflow in XACT name parsing, and a genuine data race in \cnaclass{WaveBank}'s
own cache --- \texttt{Dispose()} could race a concurrently in-flight
\texttt{GetSoundEffect()} decode --- fixed with a mutex. Two intermittent, undiagnosed
crashes in the full test suite remain open at the time of writing, both resistant to
diagnosis under both \texttt{gdb} and Valgrind so far.

\section{Microphone}

\cnaclass{Microphone} provides real SDL3 capture-device enumeration and streaming, a
\texttt{BufferReady} event, and enforces a buffer-duration constraint of 100 to 1000
milliseconds in multiples of 10 milliseconds --- matching real XNA's own documented
constraint on this API exactly.
